\documentclass[10pt,a4paper]{amsart}
\usepackage[margin=19mm]{geometry}
\usepackage{amsmath,amssymb,mathtools}
\usepackage[scr=rsfs,cal=euler]{mathalfa}
\usepackage{xfrac}
\usepackage[unicode,hidelinks,bookmarks=false]{hyperref}
\newcommand{\ie}{i.e.\ }
\newcommand{\cf}{cf.\ }
\newcommand{\resp}{respectively\ }
\newcommand{\Subsection}{\subsection}
\providecommand{\nobreakdash}{\nobreak\hbox{-}}
\setlength{\emergencystretch}{2em}
\multlinegap0pt
\usepackage{mathtools}
\usepackage{bbm}
\usepackage{amssymb}
\newtheorem{proposition}{Proposition}
\newtheorem{definition}{Definition}
\newtheorem{remark}{Remark}
\DeclareMathOperator{\Cl}{Cl}
\DeclareMathOperator{\Pic}{Pic}
\newcommand{\Z}{\mathbb{Z}}
\DeclareMathOperator{\lcm}{lcm}
\title{Calabi-Yau complete intersections in \\ fake weighted projective spaces}
\author{Marco Ghirlanda}
\date{Source: arXiv:2602.13057v1 (CC BY 4.0)}
\begin{document}
\maketitle

\noindent\emph{Source and licence.} Calabi-Yau complete intersections in \\ fake weighted projective spaces. Version arXiv:2602.13057v1 (CC BY 4.0), distributed by arXiv under the Creative Commons Attribution 4.0 International licence, http://creativecommons.org/licenses/by/4.0/, as recorded in the arXiv licence metadata for this exact version and shown on the abstract page https://arxiv.org/abs/2602.13057v1. Full licence notice: \texttt{licenses/arxiv-2602.13057-CC-BY-4.0.txt}.\par\smallskip

\noindent\textit{Editorial scope.} Counted unit: the article's proposition codimension (source0.tex lines 338--347). The article's complete original proof of it is delivered with it (source0.tex lines 349--366, one \texttt{proof} environment of 18 lines). No mathematics is written, repaired or re-derived: the statement and the proof are the article's own lines, in the article's own notation, and no retained line is substituted or re-flowed.  Retained uncounted as the source-local context the proof needs: lines 168--183; lines 187--191; lines 193--201; lines 208--219.  Nothing was omitted from any retained span, so the package carries no omission record.  No retained line contains a citation, so the wrapper carries no bibliography and no bibliography entry.  No retained line contains a figure environment, an image-inclusion command or drawing code, so no image is delivered and the package declares no asset dependency.  Editorial formatting only, in addition: an AMS \texttt{amsart} preamble that restates the article's own declarations and macros used by the retained lines, namely the environments \texttt{proposition}, \texttt{definition}, \texttt{remark} on separate counters, together with the standard packages those lines need.  The retained environments print in this document as Proposition 1, Definition 1, Remark 1 in order of appearance, whereas the article prints the same items with the numbering of its own full document.  The complete native source of the article as distributed in the arXiv e-print for this exact version is bundled unchanged as \texttt{sources/194607/source0.tex}.
\begin{proposition} \label{prop:pic}
    Let $Z$ be a fwps with $\Cl(Z)=\Z\oplus \Z/\mu_1\Z\oplus\dots\oplus\Z/\mu_r\Z$ and degree matrix $Q=\begin{bmatrix}
        \omega_0 & \dots & \omega_n
    \end{bmatrix}$, where $\omega_i=(w_i,\eta_i)$.
    Let $0\leq\eta_{ij}'<\mu_j$ be the integer representing $\eta_{ij}$ and set
    \[
    \begin{array}{lcll}  
    L & \coloneqq & \lcm(w_0,\dots,w_n),
    \\[6pt]
    M_i & \coloneqq & \dfrac{\mu_i}{\gcd\left(\mu_i,\dfrac{L}{w_0}\eta_{0j}',\dots,\dfrac{L}{w_n}\eta_{nj}'\right)}, \quad 1 \le i \le r,
    \\[24pt]
    M & \coloneqq & \lcm(M_1,\dots,M_r).
    \end{array}
    \] 
    Then the Picard group of $Z$ is the subgroup of $\Cl(Z)$ generated by the element~$(LM,0)$.
\end{proposition}

\begin{definition}\label{def:nef-part}
    Let $Z$ be a fwps with degree matrix $Q=[\omega_0,\dots,\omega_n]$. A \emph{nef-partition} of $Z$ is a partition $p_1\sqcup\dots\sqcup p_s=\{0,\dots,n\}$ such that
    $$\omega_{p_j}\coloneq\sum_{i\in p_j}\omega_i\in\Pic(Z), \hspace{5pt} j=1,\dots,s.$$
    The \emph{Calabi-Yau complete intersection (CYCI)} arising from such a partition is the family of complete intersections of general sections of $\omega_{p_1},\dots,\omega_{p_s}$ in $Z$.
\end{definition}

\begin{remark}\label{rem:linear-elimination} 
If $\omega_{p_j}=\omega_i$ for some indices $i$ and $j$, then a general section of~$\omega_{p_j}$ contains a monomial supported on the single
variable corresponding to $\omega_i$.
One can then eliminate this variable and the corresponding equation, obtaining an
isomorphic Calabi--Yau complete intersection in a fwps of dimension $d-1$ and codimension
$s-1$. Therefore, going forward we will restrict Definition \ref{def:nef-part} to nef-partitions for which
$\omega_{p_j}\neq \omega_i$ for all $j$ and all $i$, excluding in particular any partition having a
block of size one.
\end{remark}

\begin{proposition}\label{prop:nef-part}
    Let $Z$ be a fwps with $\Cl(Z)=\Z\oplus \Z/\mu_1\Z\oplus\dots\oplus\Z/\mu_r\Z$ and degree matrix $Q=\begin{bmatrix}
        \omega_0 & \dots & \omega_n
    \end{bmatrix}$. Then $p_1\sqcup\dots\sqcup p_s=\{0,\dots,n\}$ is a nef-partition if and only if, with the notation of Proposition \ref{prop:pic}, we have
    $$
    \begin{cases}
        L & | \ \ \sum_{i\in p_j} w_i,\quad j=1,\dots,s,\\[3pt]
        M_k & | \ \ \dfrac{\sum_{i\in p_j}w_i}{L},\quad j=1,\dots,s,\ k=1,\dots,r,\\[5pt]
        \sum_{i\in p_j}\eta_{ik} & =\ \ 0\in\Z/\mu_k\Z,\quad j=1,\dots,s,\ k=1,\dots,r.
    \end{cases}
    $$
\end{proposition}

\begin{proposition}\label{prop:maximal codimension}
    Let $Z$ be a fwps of dimension $n$ with $\Cl(Z)=\Z\oplus \Z/\mu_1\Z\oplus\dots\oplus\Z/\mu_r\Z$ admitting a nef-partition $p_1\sqcup\dots\sqcup p_s=\{0,\dots,n\}$. Then $2s\leq n+1$. Furthermore, a fwps $Z$ of dimension $2s-1$ admits such a nef partition if and only if $\mu_1=\dots=\mu_r=2$ and $Z$ admits a degree matrix of the form
    $$
    Q = \begin{bmatrix}
    \mathbbm{1}_{s} & \mathbbm{1}_s \\
    A & A
    \end{bmatrix},\quad \mathbbm{1}_s=[1,\dots,1],
    $$
    where $A$ is an $r\times s$ matrix in $\Z/2\Z$ whose columns affinely span the whole $(\Z/2\Z)^r$.
\end{proposition}

\begin{proof}
    First, notice that, if $2s>n+1$, at least one of the blocks of the nef-partition would have size one, contradicting Remark \ref{rem:linear-elimination}.
    
    Now assume that $n=2s-1$. Then each block of the nef-partition has size two.
    Let $\omega_i=(w_i,\eta_i)$ and $\omega_j=(w_j,\eta_j)$ be two weights forming a block.
    By Proposition \ref{prop:nef-part}, we have $L \mid w_i + w_j$.
    By definition we have $w_i$, $w_j \le L$. Therefore, either~$w_i = w_j = L$ or $w_i+w_j=L$. In the latter case we have $w_i \mid L-w_i=w_j$ and viceversa, implying $w_i=w_j=\frac{L}{2}$.
    If all pairs $(w_i,w_j)$ were of the type $\left(\frac{L}{2},\frac{L}{2}\right)$, the lowest common multiple of the weights would also be $\frac{L}{2}$, contradicting the definition of $L$. Hence, at least one pair is of the type $(L,L)$. Also, since $\gcd(w_0,\dots,w_s)=1$, we have that either $L=1$ or $L=2$ and at least a pair has the form $(1,1)$. In the latter case, the sum of the pair would coincide with one of the other $w_i$, contradicting Remark \ref{rem:linear-elimination}. Hence, we conclude that $w_i=1$ for all $i$.

    By Proposition \ref{prop:nef-part}, it follows that $\mu_i=2$ for all $i$, and each block of the nef-partition is of the form $(\omega_i,\omega_i)$ with $\omega_i=(1,\eta_i)$. After appropriately permuting the columns of the degree matrix, we can therefore reach the form
    $$
    Q = \begin{bmatrix}
    \mathbbm{1}_{s} & \mathbbm{1}_s \\
    A & A
    \end{bmatrix}.
    $$
    We conclude by noticing that the almost freeness of $Q$ is equivalent to the condition that the columns of $A$ span the whole $(\Z/2\Z)^r$.
\end{proof}

\end{document}
